
\begin{figure*}[tbp]
    \centering
    \includegraphics[width=\textwidth]{stability_measures.png}
    \caption{Illustrations of some common stability measurements in computational material science. (a) The formation energy $E_f$ of a compound is defined by referencing to the energies of its constituent elements in their standard form, which is demonstrated here for the chemical reaction of Si$(s)$+O$_{2}(g)\to$SiO$_{2}(s)$ that corresponds to  the formation of silicon dioxide. (b) (Top) Energy landscape for the intermixing between cubic halide perovskites of \ce{CsPbI3} and \ce{CsPbBr3} forming the solid solution of CsPb(I$_{1-x}$Br$_{x}$)$_3$.\cite{yang2020correlating} The red dots highlight the convex hull of the energy landscape. (Bottom) Exemplary atomistic structures of the  CsPb(I$_{1-x}$Br$_{x}$)$_3$ solid solutions with different concentrations ($x$) and stabilities. (c) Thermodynamic stability diagram as a function of O and La chemical potentials built for \ce{LaMnO3} (001) surfaces.\cite{piskunov2008electronic} It contains comparison of stability of both LaO-- and \ce{MnO2}--terminated (001) surfaces and accounts for precipitation conditions for La and Mn metals and their trivalent oxides \ce{La2O3} and \ce{Mn2O3}. Stable region is shown as the hatched area between \ce{La2O3} and \ce{Mn2O3} precipitation lines. The right side shows a family of oxygen chemical potentials under different conditions. The label $m$ indicates the \ce{O2} gas partial pressure: $10^{m}$ atm. Red gray line corresponds to oxygen partial pressure $p=0.2p_0$ as in the ambient atmosphere. (d) Range of the chemical potentials $\mu_{\mbox{\scriptsize{Sn}}}$ and $\mu_{\mbox{\scriptsize{Ba}}}$ (yellow shaded region)  in which the transparent oxide \ce{BaSnO3}  \cite{scanlon2013defect} is chemically stable. The verticies of the triangle are determined by the formation enthalpy $\Delta H_f$ of \ce{BaSnO3}. (e) Formation energies of the native defects in \ce{BaSnO3} as a function of the Fermi level at two diffent chemical potentials A and B as highlighted in (d).}
    \label{fig:stability_measures}
\end{figure*}
 
Practically, the potential energy of a material $E(\bm{R})$  can be computationally determined from its atomic configurations based on different level of approximations, which will be  discussed with details in \cref{chapter:energy_models}. However, regardless of what approximation is adopted,  the absolute value of the potential energy $E$ itself that is calculated for a given atomic configuration $\{\bm{R}\}$ bears little physical significance, and must be interpreted by comparing it to the energy of a reference state in order to provide a meaningful measure of the material stabilities. Conceptually, this is equivalent to the consideration of a  reaction $\mathcal{A}\to\mathcal{B}$, where the reaction energy is always measured as the difference in energy of the product(s) with respect to the reactant(s) as $\Delta E=E[\bm{R}(\mathcal{A})]-E[\bm{R}(\mathcal{B})]$. Although practically there may not be a chemical reaction happening, for example, in the case where $\mathcal{A}$  and  $\mathcal{B}$  corresponds to two different conformers of a molecules  or polymorphs of a crystal.  In this section, some important energetic  merits for material stability that frequently arises in computational modellings will be briefly reviewed, and we focus our discussions on the energetic measures at 0 K. 

The \textit{formation energy} ($E_f$) of a compound, for instance, A$_l$B$_m$C$_n$, is defined for the formation reaction from its constituent elements in their standard form, i.e.
\begin{equation}\label{eq:formation_chemical_reaction}
    l\mbox{A}+m\mbox{B}+n\mbox{C}\to \mbox{A}_l\mbox{B}_m\mbox{C}_n,
\end{equation}
with 
\begin{equation}\label{eq:formation_energy}
    E_{f}\left(\mbox{A}_l\mbox{B}_m\mbox{C}_n\right)=\frac{1}{N}\left[E\left(\mbox{A}_l\mbox{B}_m\mbox{C}_n\right)-lE(\mbox{A})-mE(\mbox{B})-nE(\mbox{C})\right],
\end{equation}
where $N$ is the total number of atoms in the (simulation) cell. This is demonstrated in \cref{fig:stability_measures}(a) for \ce{SiO2}.  Compounds that possess $E_f<0$ are generally considered to be stable (against decomposition to its constituent elements). For crystalline compounds, $E_f$ is sometimes also referred as the \textit{cohesive energy} or the \textit{lattice energy}. If all the reactants on the left--hand--side of the reactions are further constrained to the gas phase of individual element, then $E_f=-\Delta H_{sub}$ where $\Delta H_{sub}$ is the sublimation enthalpy of the solid. 

Note that although we have chosen to illustrate the concept of $E_f$ using a compound with strong chemical bond (\ce{SiO2} in this case), it is equally applicable to aggregates bonded together with weak bonding, such as van der Waals (vdW) interactions. Normally, this is referred to as the \emph{interaction} or \emph{binding energy} $E_{b}$, such as the binding energies for a graphene bilayer, water dimer, or the cohesive energy of molecular crystals. The vdW binding energy is typically an order or two weaker than a strong chemical bond.  For example, the bonding energy of a C--C single bond is around 90 kcal/mol compared to $\sim$2.5--3.0 kcal/mol for the binding energy of a benzene dimer.

A more generalised version of the reaction (\cref{eq:formation_chemical_reaction}) to be considered, that is better aligned with the concept of phase diagram [\cref{fig:energy_landscape_examples}(f)] in material science, is to rewrite reaction (\cref{eq:formation_chemical_reaction}) as 
\begin{equation}\label{eq:convex_hull_reaction}
n_{1}\mathcal{P}_{1}+n_{2}\mathcal{P}_2+\cdots\to \mathcal{C}^{(i)},
\end{equation}
which represents the formation of a new compound $\mathcal{C}$ in a particular phase $(i)$, or, the $i$--th configuration of a solid solution, with a fixed chemical composition, from the stable parent phases  $\{\mathcal{P}_{1},\mathcal{P}_{2},\ldots\}$ on the phase diagram. In this case, $\mathcal{P}$ could be either an element or a compound, and $\mathcal{C}$ could occur in multiple different phases or atomic configurations. For example, \cref{fig:stability_measures} (b) demonstrates the formation of  mixed halide solid solutions, where  $\mathcal{C}$ in (\cref{eq:convex_hull_reaction}) corresponds to CsPb(I$_{1-x}$Br$_{x}$)$_3$ with a fixed halide concentration $x$, while $\mathcal{P}_{1}$ and $\mathcal{P}_{2}$ are the parent \ce{CsPbBr3} and \ce{CsPbI3} perovskites. With this in mind, a negative (positive) reaction energy, or most commonly referred to as the \emph{mixing enthalpy} 
\begin{equation}\label{eq:mixing_enthalpy}
\Delta H_{mix}^{(i)}=E[\mathcal{C}^{(i)}]-\sum_k n_k E[\mathcal{P}_{k}]
\end{equation}
indicates whether the compound is less (more) likely to undergo chemical decompositions into its parent phases.   By further sampling the possible structures $\{\mathcal{C}^{(1)},\mathcal{C}^{(2)},\mathcal{C}^{(3)},\ldots\}$ at a particular chemical composition to obtain their corresponding mixing enthalpies $\{\Delta H_{mix}\}=\{\Delta H_{mix}^{(1)},\Delta H_{mix}^{(2)},\ldots\}$, the position for the \textit{convex hull} of stability can be identified as the point corresponding to $\min\{\Delta H_{mix}\}$ [highlighted with red dots in \cref{fig:stability_measures}(b)], whereas
\begin{equation}\label{eq:convex_hull_energy}
\Delta E_{hull}^{(i)}=\Delta H_{mix}^{(i)} - \min\{\Delta H_{mix}\}
\end{equation}
is the convex hull energy for a given material, which is typically applied as a measure of the material's metastability in its $i$--th phase.\cite{sun2016thermodynamic}

Beyond the phase stabilities mentioned above, for some materials that are particularly prone to chemical degredations under the ambient conditions, it is also necessary to further compare their phase stabilities with respect to the energies of the degradation reactions for these materials, in order to obtain a more complete assessment on their chemical stabilities. Considering this effect has proven to be critical for designing materials with outstanding physical properties that have been intensively studied in the recent years, including the oxidative degradation of halide perovskites as photovoltaic materials,\cite{senocrate2019thermochemical,senocrate2018interaction} black phosphorous,\cite{zhou2016light,zhang2018degradation} as well as thermoelectric SnSe.\cite{li2016high} 

Instead of fixing the energies of the  elemental reference states in \cref{eq:formation_energy}, we recognise that these energies are thermodynamic state variables, where the most commonly used one being the the \emph{chemical potential} $\mu$. Practically, this means that rather than being a single--valued quantity, the formation energy $E_f$ is a multi--dimensional function of the chemical potentials of its constituent elements, \emph{e.g.}:
\begin{equation}\label{eq:formation_energy_chem_pot_func}
E_f(\mbox{A}_l\mbox{B}_m\mbox{C}_n; \mu_A,\mu_B,\mu_C) =\frac{1}{N}\left[ E\left(\mbox{A}_l\mbox{B}_m\mbox{C}_n\right)  -l\mu_{A}-m\mu_B-n\mu_C\right],
\end{equation}
whereby the stability region for a specific structure is the range of $\mu$'s within which $E_f<0$. It is apparent from \cref{eq:formation_energy_chem_pot_func} that  different compounds/phases formed with a fixed set of elements  will have different energies $E$, as such the stability regions for these different compounds/phases will cover different ranges of accessible $\mu$'s on the phase diagram.  This further means that for a given structure to be \emph{the} most stable one, its $E_f$ must be more negative than those for all other competing structures at some $\mu$'s, which may further narrow down the stability region for a given structure. Therefore, the ability to map out the the stability regions for different competing structures on the phase diagram  spanned by the chemical potentials of the constituent elements can provide the important guideline in determining the suitable experimental conditions  for fabricating a specific material/structure. An example of this is a (computationally determined) stability diagram for the (001) surfaces of a complex oxide \ce{LaMnO3} [\cref{fig:stability_measures}(c)],\cite{piskunov2008electronic} showing two distinct stability regions corresponding to the  surfaces with either \ce{LaO}-- and \ce{MnO2}--terminations, respectively, which can be applied to guide the experimental synthesis of complex oxide thin films with vacuum depositions. In particular, as shown in \cref{fig:stability_measures}(c), the chemical potentials of the oxygen gas can be determined by its partial pressure and temperature. 

Finally, we briefly consider the \textit{formation energies of charged defects}  in crystalline solids. The presence of point defects in materials plays an important role in determining its physical properties. The formation of defects in the materials is generally unavoidable, but may be controlled with a good understanding in their formation thermodynamics. Given its profound importance, there already exist many prior reviews specifically devoted to the discussions of defect chemistry and physics in crystalline solids.\cite{freysoldt2014first,anand2022convergent} As such, we will only briefly highlight the key ideas behind the defect thermodynamics here. Following Annand \emph{et al.}\cite{anand2022convergent}, consider the formation reaction of an oxygen vacancy V$_{\mbox{\footnotesize{O}}}$ in a binary oxide that are both mass and charge conserved:
\begin{equation*}
    \mbox{MO}\to\mbox{MO}_{1-\delta}+\delta \mbox{O} + q\delta \mbox{e}^{-},
\end{equation*}
where $q$ is the charge transferred from the solid into the reservoir upon the formation of the charged defect, and in this case $q=2$ Now if one follows the formulation of formation energy as in \cref{eq:formation_energy_chem_pot_func}, the  formation energy of a V$_{\mbox{\footnotesize{O}}}$ in the binary oxide MO would be
\begin{equation}\label{eq:formation_defect_chem_pot}
E_f\left(\mbox{V}_{\mbox{\footnotesize{O}}}\right)=\left[E\left(\mbox{MO}_{1-\delta}\right)-E\left(\mbox{MO}\right)\right]-\delta \mu_{\mbox{\footnotesize{O}}}-q\delta\mu_{e}.
\end{equation}
Here, $E\left(\mbox{MO}_{1-\delta}\right)$ and $E\left(\mbox{MO}\right)$ are the total energies of the solid with and without the oxygen vacancy. Generally, the chemical potential of an electron in a semiconductor is defined as the Fermi level $(E_F)$ with respect to the energy position of the valence band maximum (VBM), hence $\mu_e=E_{F}+E_{\mbox{\footnotesize{VBM}}}$. Replacing this definition of $\mu_e$ into \cref{eq:formation_defect_chem_pot}, it is apparent that $E_f\left(\mbox{V}_{\mbox{\footnotesize{O}}}\right)$ will not only depend on the chemical potential of the oxygen that is being removed from the solid, but also the Fermi level and the charge state of the oxygen vacancy. This concept is illustrated in \cref{fig:stability_measures}(d--e) where the formation energy landscapes of the native defects (with respect to $E_{F}$ and $q$) in perovskite \ce{BaSnO3} are compared at two different sets of chemical potentials at which the \ce{BaSnO3} phase is stable. It demonstrates that the native defects in undoped  \ce{BaSnO3} are energetically stable when $E_{F}$ is close to the VBM, posing a challenge to achieve $n$--type electrical conductivity in undoped  \ce{BaSnO3}.

